\chapter{Talk to the project agent}\label{ch:chat}

\emph{You want answers about your code, and you want to turn an idea into a well-written issue.}
The project's chat agent is Claude Code working in your own checkout of the repository, with your
Claude credential. It can read the code, run commands, and use Agern itself: list and create
issues, start, approve and cancel runs.

\section{Ask about the code}\label{sec:ask}

\begin{enumerate}
  \item Open a project. If you have no chat there yet, choose \ui{Start chat}; \ui{+ New chat}
    starts another, fresh conversation later.
  \item Type your question into the prompt at the bottom (\ui{Ask the agent, or /factory \ldots})
    and press Enter. Shift+Enter makes a new line.
  \item Watch the answer stream in: the agent's thinking, each tool it uses (reading files,
    searching, running commands) with its result, and the reply. The line under the reply shows
    how long it took, what it cost and how many turns it took (figure~\ref{fig:check-chat} in
    chapter~\ref{ch:accounts}).
  \item Press Esc to interrupt the agent if it heads the wrong way, and tell it what you meant.
\end{enumerate}

A chat keeps its context: follow-up questions build on what came before. An idle chat's process
is stopped after a while; your next message resumes the conversation. Others can read your chat,
but only you can write in it.

\section{Attach screenshots and files}\label{sec:attachments}

\begin{enumerate}
  \item Choose \ui{+} next to the prompt (\ui{Attach files}), or paste or drop files onto the chat.
  \item The files appear as chips above the prompt while they upload. Write your message and send it.
  \item Images are shown to the agent as images (very large ones as a smaller copy); other files are
    saved for it to read (figure~\ref{fig:chat-attachment}).
\end{enumerate}

By default a file may have up to 20~MB, a message up to 50~MB in up to ten files. Commands
starting with \cmd{/factory} take no attachments.

\screenshot[0.8\linewidth]{30-chat-attachment-desktop}{A message with an attached screenshot, and the agent's answer.}{chat-attachment}

\section{Turn an idea into an issue}\label{sec:new-issue}

There are two ways.

\begin{enumerate}
  \item \textbf{With the agent.} Describe the idea in the chat and ask the agent to file an issue,
    for example: ``Write an issue for sharing a wishlist by link. Look at how carts are stored
    first.'' The agent reads the code, drafts the issue and creates it on GitHub with your token;
    it appears in the issue list at once. This usually gives the better issue, because the agent
    knows the code.
  \item \textbf{By hand.} Choose \ui{+ New} above the issue list. In the dialog
    (figure~\ref{fig:new-issue}) choose the kind (\ui{feature}, \ui{bug} or \ui{task}), enter a
    \ui{Title} and a \ui{Description} in Markdown, and choose \ui{Create issue}.
\end{enumerate}

A good issue says what is wrong or wanted and how you will know it is done; the spec-writer of a
run turns it into a full specification (chapter~\ref{ch:run}).

\screenshot[0.85\linewidth]{23-new-issue-desktop}{Creating an issue by hand.}{new-issue}

\section{Factory commands}\label{sec:factory-commands}

The chat also understands a few commands that act directly, without asking the model. Type
\cmd{/} at the start of the prompt to see them (figure~\ref{fig:autocomplete}):

\begin{center}
\begin{tabular}{@{}ll@{}}
  \toprule
  Command & What it does\\
  \midrule
  \cmd{/factory run <issue>} & Start a pipeline run for an issue\\
  \cmd{/factory status} & List runs in this project\\
  \cmd{/factory approve [runId]} & Approve a run waiting for a human\\
  \cmd{/factory cancel <runId>} & Cancel a run\\
  \cmd{/factory help} & Show factory commands\\
  \bottomrule
\end{tabular}
\end{center}

\screenshot{10-project-autocomplete-desktop}{The factory commands, suggested as you type.}{autocomplete}

\whatyouhave
You can ask the agent about the code, show it screenshots and files, and turn an idea into an
issue, by asking the agent or with \ui{+ New}. Next, hand that issue to the pipeline:
chapter~\ref{ch:run}.
